% ID: SEL-D0055-OA-E03
% ESTADO: cotejada. Segunda pasada visual y borrador compilado revisados.
% Fuente: Selectividad/sel_2011_matematicas_aplicadas/122A-EXAMEN 6.pdf
% SHA-256: aebcbfddd4d5db5f5a10b4662e9e12b2ee362bbf574b751bccd6d78594766ad3
% Página: 1; opción A; ejercicio original 3.

Un examen consta de una parte teórica y una parte práctica. La probabilidad de que se apruebe la parte teórica es 0.7 y la de que se apruebe la parte práctica 0.75. Se sabe que el 50\% de los alumnos ha aprobado ambas.
\begin{enumerate}[label=\alph*)]
\item \textbf{(0.75 puntos)} Calcule la probabilidad de aprobar alguna de las dos partes.
\item \textbf{(0.75 puntos)} Calcule la probabilidad de aprobar la parte práctica sabiendo que no se ha aprobado la parte teórica.
\item \textbf{(1 punto)} \textquestiondown{}Son independientes los sucesos ``aprobar parte teórica'' y ``aprobar parte práctica''?
\end{enumerate}
